\problemname{Atlantic Email}

\illustration{0.38}{erste-email.jpg}{
    CC BY-SA 4.0 by Reinhard Ferdinand on \href{https://commons.wikimedia.org/wiki/File:Karlsruhe_erste_E-Mail_Deutschlands.JPG}{Wikimedia Commons}
}

% optionally define variables/limits for this problem
\newcommand{\maxn}{10\,000}
The year is 1984. Your boss, Professor Werner Zorn, has been pushing for establishing connections to international networks.
His newest project is for Karlsruhe University\footnote{Former university in Karlsruhe that later merged to become Karlsruhe Institute of Technology.}
to receive the first email from the US Computer Science Network (CSNet).

The email will be sent by Laura Breeden (Massachusetts Institute of Technology) to the Karlsruhe University computers.
For simplicity, the email only consists of $n$ lowercase letters. It should be encoded
so it can be sent across $5$ transatlantic binary channels.

You want the communication to be done efficiently, so you require that across each channel,
no more than $n+10$ bits are sent.
The channels are working flawlessly and losslessly, so for each of the $5$
channels, the bits will arrive intact in Karlsruhe.

Everything was almost ready, when you noticed that the wires coming into the
computer in Karlsruhe are a mess, and so you do not know which of the $5$ channels is which.
Thus, as the receiver, you will receive the binary strings of the $5$ channels in arbitrary order.
Instead of untangling the mess, you think it is a better idea to just fix it in software.
Design a communication protocol that can transfer the email correctly and efficiently.

\begin{Input}
    The input consists of:
    \begin{itemize}
        \item One line with the action that your program needs to perform:
        either the string ``\texttt{send}'' or ``\texttt{receive}'',
        which denotes whether you are at the sending or receiving end.

        \item If the action is ``\texttt{send}'':
        \begin{itemize}
        \item One line with one integer $n$ ($1\leq n \leq \maxn$), the length of the email.
        \item One line with a string $s$ of length $n$, the email to send. The
        email only consists of English lowercase letters (\texttt{a-z}).
        \end{itemize}
        \item If the action is ``\texttt{receive}'':
        \begin{itemize}
        \item Five lines, each with a string $b$ ($1 \leq |b| \leq n+10$), only consisting of digits $0$ and $1$, the five binary strings you used to encode the email, in an arbitrary order.

        \end{itemize}

    \end{itemize}
    Note that when receiving, you have to infer $n$ from the received binary
    strings, as this is not known to you up front.

    This is a multi-pass problem. Your program will be invoked multiple times,
    possibly more than twice.
    It is guaranteed that the first pass is a ``\texttt{send}'' action,
    and that each subsequent pass is a ``\texttt{receive}'' action.

    For testing purposes, the number and input of subsequent passes will depend
    on the output of your submission.

    A testing tool is provided to help you develop your solution.
\end{Input}

\begin{Output}
    If the action is ``\texttt{send}'', output $5$ bit strings $b_1, \dots, b_5$
    ($1 \leq |b_i| \leq n+10$ for each $i$), only consisting of digits $0$ and $1$, the bit strings to send.

    If the action is ``\texttt{receive}'', output the original email $s$, a
    string only consisting of English lowercase letters (\texttt{a-z}).
\end{Output}

